半径为 \(a>0\) 的理想导体实心球置于真空，球心在原点；球外一正点电荷 \(q>0\) 放在 \(z=b>a\) 的轴线上。忽略辐射及导体材料微观效应，静电平衡时外部电势满足球面的等势条件。考察两种情形：A 为导体球接地，B 为导体球与外界绝缘且初始净电荷为零。
(1) 对接地球构造 Kelvin 镜像的位置与电荷量，并证明球面电势严格为零；(2) 求接地球诱导的总电荷、导体对点电荷的力及相互作用势能；(3) 对孤立中性球给出额外镜像构造，求外部电势、导体对点电荷的力，比较它与接地球的吸引强弱；(4) 讨论 \(b\gg a\) 与 \(b\to a^+\) 极限，并解释不能将镜像电荷作为真实可自由运动的点电荷处理。
